\section{Background and Preliminaries}
\label{sec:background}

Below is provided a self-contained review of the mathematical foundations of group-equivariant
convolutional networks. Readers familiar with \citet{cohen2016group} may skip to
Section~\ref{sec:method}.

\subsection{Groups and Group Actions}

A \emph{group} \((G, \cdot)\) is a set \(G\) equipped with a binary operation \(\cdot\) that is
associative, has an identity element \(e \in G\), and for which every element \(g \in G\)
has an inverse \(g^{-1} \in G\). Groups formalise the notion of \emph{symmetry}: the elements
of \(G\) are the symmetry transformations, and the group operation corresponds to their
composition.

A group \emph{acts} on a set \(X\) if there is a map \(\rho: G \times X \to X\),
written \((g, x) \mapsto g \cdot x\), that satisfies \(e \cdot x = x\) and
\(g \cdot (h \cdot x) = (gh) \cdot x\) for all \(g, h \in G\) and \(x \in X\).

\paragraph{The translation group \(\Z^2\).}
Standard 2-D convolutions are defined over functions \(f: \Z^2 \to \R^K\) and exploit the
action of the translation group \(\Z^2\) on pixel coordinates. The group operation is
vector addition \((m, n) + (p, q) = (m+p, n+q)\).

\paragraph{The group \(p4\).}
The wallpaper group \(p4\) consists of all compositions of integer translations and rotations
by multiples of \(90^\circ\) about any lattice point. Each element can be parameterised as
a triple \((r, u, v)\) with rotation index \(r \in \{0, 1, 2, 3\}\) and translation
\((u, v) \in \Z^2\), corresponding to the transformation matrix
\begin{equation}
  g(r, u, v) =
  \begin{pmatrix}
    \cos\!\tfrac{r\pi}{2} & -\sin\!\tfrac{r\pi}{2} & u \\
    \sin\!\tfrac{r\pi}{2} &  \cos\!\tfrac{r\pi}{2} & v \\
    0                     &  0                     & 1
  \end{pmatrix}.
  \label{eq:p4-matrix}
\end{equation}
The group operation is matrix multiplication. The group \(p4\) has index 4 over
\(\Z^2\) (four rotational orientations per translation).

\paragraph{The group \(p4m\).}
The group \(p4m\) extends \(p4\) by including mirror reflections in addition to rotations,
yielding an index-8 extension of \(\Z^2\). Each element is parameterised as
\((m, r, u, v)\) with flip bit \(m \in \{0, 1\}\):
\begin{equation}
  g(m, r, u, v) =
  \begin{pmatrix}
    (-1)^m \cos\!\tfrac{r\pi}{2} & -(-1)^m \sin\!\tfrac{r\pi}{2} & u \\
    \sin\!\tfrac{r\pi}{2}        &  \cos\!\tfrac{r\pi}{2}         & v \\
    0                            &  0                             & 1
  \end{pmatrix}.
  \label{eq:p4m-matrix}
\end{equation}
A \(p4m\)-equivariant feature map therefore carries 8 channels per ``logical'' feature,
corresponding to 4 rotations \(\times\) 2 reflection states.

\subsection{Equivariance}

Let \(\Phi: \mathcal{X} \to \mathcal{Y}\) be a learned map between function spaces. We say
\(\Phi\) is \emph{equivariant} to the group action of \(G\) if there exist representation
operators \(T_g\) on \(\mathcal{X}\) and \(T'_g\) on \(\mathcal{Y}\) such that
\begin{equation}
  \Phi(T_g\, x) = T'_g\, \Phi(x) \quad \forall\, g \in G,\; x \in \mathcal{X}.
  \label{eq:equivariance}
\end{equation}
Intuitively, transforming the input and then passing it through \(\Phi\) gives the same result
as passing it through \(\Phi\) first and then applying the corresponding output transformation.
\emph{Invariance} is the special case \(T'_g = \mathrm{Id}\) for all \(g\).

Equivariance is more useful than invariance in intermediate layers because it preserves
spatial information: the network retains the \emph{pose} of detected features, allowing
subsequent layers to reason about their relative configuration. Invariance to the full group
is typically enforced only at the final pooling stage.

\subsection{Standard Convolution}

For an input feature map \(f: \Z^2 \to \R^K\) and a filter \(\psi: \Z^2 \to \R^K\),
the standard \emph{cross-correlation} (commonly called convolution in the deep-learning
literature) is
\begin{equation}
  [f \star \psi](x) = \sum_{y \in \Z^2} \sum_{k=1}^{K} f_k(y)\,\psi_k(y - x).
  \label{eq:standard-conv}
\end{equation}
This operation is equivariant to translations in \(\Z^2\), because shifting the input
\(f\) by \(t\) shifts the output by \(t\) as well. However, it is not equivariant to
rotations: rotating the input generally produces a different output than rotating the output
of the unrotated input.

\subsection{Group Convolution}

The \gcnn{} framework \citep{cohen2016group} generalises Eq.~\eqref{eq:standard-conv}
to arbitrary groups \(G\). The central operation is the \emph{\(G\)-convolution} (or
\emph{group cross-correlation}).

\paragraph{First layer (\(\Z^2 \to G\) convolution).}
Given a planar input \(f: \Z^2 \to \R^K\) and a filter \(\psi: \Z^2 \to \R^K\), the
\(\Z^2 \to G\) group convolution lifts the feature map from the plane to the group:
\begin{equation}
  [f \star_G \psi](g) = \sum_{y \in \Z^2} \sum_{k=1}^{K} f_k(y)\,\psi_k(g^{-1} y),
  \qquad g \in G.
  \label{eq:lift-conv}
\end{equation}
For \(G = p4\), this computes four responses per spatial location, one for each of the four
rotations \(r \in \{0,1,2,3\}\), effectively correlating the input with four rotated copies
of the same filter.

\paragraph{Subsequent layers (\(G \to G\) convolution).}
In deeper layers, both the feature map and the filter are defined on \(G\). The
\(G \to G\) group convolution is
\begin{equation}
  [f \star_G \psi](g) = \sum_{h \in G} \sum_{k=1}^{K} f_k(h)\,\psi_k(g^{-1} h),
  \qquad g \in G.
  \label{eq:group-conv}
\end{equation}
This can be interpreted as a convolution on the group manifold: the filter \(\psi\) is
shifted across \(G\) by left multiplication.

\paragraph{Equivariance guarantee.}
One can show \citep{cohen2016group} that the \(G \to G\) group convolution
(Eq.~\ref{eq:group-conv}) satisfies Eq.~\eqref{eq:equivariance} with
\(T_g\, f = f(g^{-1} \cdot)\), the regular representation of \(G\). This guarantee holds
for arbitrary group elements, not just those seen during training, and extends to composed
architectures of arbitrary depth. Non-linearities, batch normalisation, and pooling
operations applied element-wise on \(G\)-feature maps are all equivariant, so any
sequence of such layers preserves the equivariance of the whole network.

\paragraph{Group-invariant output.}
For a classification task, the final representation must be invariant rather than
equivariant. This is achieved by a \emph{group pooling} operation that averages (or
max-pools) over the group dimension:
\begin{equation}
  \bar{f}(x) = \frac{1}{|G / \Z^2|} \sum_{r} f(g(r, x)),
  \label{eq:group-pool}
\end{equation}
collapsing the orientation axis and producing a translation-equivariant (but
rotation-\emph{invariant}) feature map suitable for global average pooling and
classification.

\paragraph{Parameter efficiency.}
Because the \(G\)-convolution shares a single filter across all \(|G / \Z^2|\) orientations,
the number of parameters for a given number of logical output channels grows by
\(|G / \Z^2|\) relative to a standard convolution. To keep total parameter counts
comparable, \citet{cohen2016group} recommend dividing the number of filters by
\(\sqrt{|G / \Z^2|}\) (e.g.\ by 2 for \(p4\), or by \(\sqrt{8}\) for \(p4m\)),
so that the product of filter count and group size remains roughly constant.
The additional expressive power then comes from the filter being evaluated at every
orientation, effectively seeing a larger variety of training patterns per parameter update.